\section{2015年高考题}
\subsection{2015年全国I卷（理）}\label{gk:2015年全国I卷（理）}


\clearpage


\subsection{2015年全国I卷（文）}\label{gk:2015年全国I卷（文）}


\clearpage


\subsection{2015年全国II卷（理）}\label{gk:2015年全国II卷（理）}


\clearpage


\subsection{2015年全国II卷（文）}\label{gk:2015年全国II卷（文）}


\clearpage


\subsection{2015年安徽卷（理）}\label{gk:2015年安徽卷（理）}


\clearpage



\subsection{2015年安徽卷（文）}\label{gk:2015年安徽卷（文）}


\clearpage


\subsection{2015年北京卷（理）}\label{gk:2015年北京卷（理）}

\begin{description}
    \item[一、]
    \begin{enumerate}
        \item 1
        \item 2
        \item 3
        \item 4
        \item 5
        \item 6
        \item 7
        \item 8
    \end{enumerate}
    \item[二、]
    \begin{enumerate}
      \addtocounter{enumi}{+8}   
        \item 1
        \item 2
        \item 3
        \item 4
        \item 1
        \item 6
    \end{enumerate}
    \item[三、]   
    \begin{enumerate}
      \addtocounter{enumi}{+14}   
        \item 1
        \item 2
        \item 3
        \item 已知函数 $f(x) = \ln \frac{1+x}{1-x}$.

\begin{enumerate}[label=(\arabic*)]
    \item 求曲线 $y = f(x)$ 在点 $(0, f(0))$ 处的切线方程；
    \item 求证：当 $x \in (0,1)$ 时，$f(x) > 2\left(x + \frac{x^{3}}{3}\right)$；
    \item 设实数 $k$ 使得 $f(x) > k\left(x + \frac{x^{3}}{3}\right)$ 对 $x \in (0,1)$ 恒成立，求 $k$ 的最大值。
\end{enumerate}
        \item 已知椭圆 $C: \frac{x^{2}}{a^{2}} + \frac{y^{2}}{b^{2}} = 1$ ($a > b > 0$) 的离心率为 $\frac{\sqrt{2}}{2}$，点 $P(0,1)$ 和点 $A(m,n)$ ($m \neq 0$) 都在椭圆 $C$ 上，直线 $PA$ 交 $x$ 轴于点 $M$.

\begin{enumerate}[label=(\arabic*)]
    \item 求椭圆 $C$ 的方程，并求点 $M$ 的坐标；（用 $m$, $n$ 表示）
    \item 设 $O$ 为原点，点 $B$ 与点 $A$ 关于 $x$ 轴对称，直线 $PB$ 交 $x$ 轴于点 $N$，问：$y$ 轴上是否存在点 $Q$，使得 $\angle OQM = \angle ONQ$？若存在，求点 $Q$ 的坐标；若不存在，说明理由。
\end{enumerate}
        \item 已知数列 $\{a_n\}$ 满足：$a_1 \in \mathbb{N}^{*}$，$a_1 \leqslant 36$，且 
\[
a_{n+1} = 
\begin{cases} 
2a_n, & a_n \leqslant 18, \\
2a_n - 36, & a_n > 18,
\end{cases} 
\quad (n = 1, 2, \dotsc).
\]
记集合 $M = \{a_n \mid n \in \mathbb{N}^{*}\}$.

\begin{enumerate}[label=(\arabic*)]
    \item 若 $a_1 = 6$，写出集合 $M$ 的所有元素；
    \item 若集合 $M$ 存在一个元素是 $3$ 的倍数，证明：$M$ 的所有元素都是 $3$ 的倍数；
    \item 求集合 $M$ 的元素个数的最大值。
\end{enumerate}
    \end{enumerate}
\end{description}
\clearpage


\subsection{2015年北京卷（文）}\label{gk:2015年北京卷（文）}

\begin{description}
    \item[一、]
    \begin{enumerate}
        \item 1
        \item 2
        \item 3
        \item 4
        \item 5
        \item 6
        \item 7
        \item 8
    \end{enumerate}
    \item[二、]
    \begin{enumerate}
      \addtocounter{enumi}{+8}   
        \item 1
        \item 2
        \item 3
        \item 4
        \item 1
        \item 6
    \end{enumerate}
    \item[三、]   
    \begin{enumerate}
      \addtocounter{enumi}{+14}   
        \item 1
        \item 2
        \item 3
        \item 
        \item  设函数 $f(x) = \frac{x^{2}}{2} - k \ln x$, $k > 0$.
\begin{enumerate}[label=(\arabic*)]
    \item 求 $f(x)$ 的单调区间和极值；
    \item 证明：若 $f(x)$ 存在零点，则 $f(x)$ 在区间 $(1, \sqrt{e}]$ 上仅有一个零点。
\end{enumerate}

        \item 已知椭圆 $C: x^{2} + 3y^{2} = 3$，过点 $D(1,0)$ 且不过点 $E(2,1)$ 的直线与椭圆 $C$ 交于 $A$, $B$ 两点，直线 $AE$ 与直线 $x=3$ 交于点 $M$.
\begin{enumerate}[label=(\arabic*)]
    \item 求椭圆 $C$ 的离心率；
    \item 若直线 $AB$ 垂直于 $x$ 轴，求直线 $BM$ 的斜率；
    \item 试判断直线 $BM$ 与直线 $DE$ 的位置关系，并说明理由。
\end{enumerate}
    \end{enumerate}
\end{description}
\clearpage


\subsection{2015年重庆卷（理）}\label{gk:2015年重庆卷（理）}


\clearpage


\subsection{2015年重庆卷（文）}\label{gk:2015年重庆卷（文）}


\clearpage


\subsection{2015年福建卷（理）}\label{gk:2015年福建卷（理）}


\clearpage


\subsection{2015年福建卷（文）}\label{gk:2015年福建卷（文）}


\clearpage


\subsection{2015年广东卷（理）}\label{gk:2015年广东卷（理）}


\clearpage


\subsection{2015年广东卷（文）}\label{gk:2015年广东卷（文）}


\clearpage


\subsection{2015年湖北卷（理）}\label{gk:2015年湖北卷（理）}


\clearpage


\subsection{2015年湖北卷（文）}\label{gk:2015年湖北卷（文）}


\clearpage


\subsection{2015年湖南卷（理）}\label{gk:2015年湖南卷（理）}


\clearpage


\subsection{2015年湖南卷（文）}\label{gk:2015年湖南卷（文）}


\clearpage


\subsection{2015年江苏卷}\label{gk:2015年江苏卷}


\clearpage


\subsection{2015年江西卷（理）}\label{gk:2015年江西卷（理）}


\clearpage


\subsection{2015年江西卷（文）}\label{gk:2015年江西卷（文）}


\clearpage


\subsection{2015年辽宁卷（理）}\label{gk:2015年辽宁卷（理）}


\clearpage


\subsection{2015年辽宁卷（文）}\label{gk:2015年辽宁卷（文）}


\clearpage


\subsection{2015年上海卷（理）}\label{gk:2015年上海卷（理）}


\clearpage


\subsection{2015年上海卷（文）}\label{gk:2015年上海卷（文）}


\clearpage


\subsection{2015年四川卷（理）}\label{gk:2015年四川卷（理）}


\clearpage


\subsection{2015年四川卷（文）}\label{gk:2015年四川卷（文）}


\clearpage


\subsection{2015年山东卷（理）}\label{gk:2015年山东卷（理）}


\clearpage


\subsection{2015年山东卷（文）}\label{gk:2015年山东卷（文）}


\clearpage


\subsection{2015年陕西卷（理）}\label{gk:2015年陕西卷（理）}


\clearpage


\subsection{2015年陕西卷（文）}\label{gk:2015年陕西卷（文）}


\clearpage


\subsection{2015年天津卷（理）}\label{gk:2015年天津卷（理）}


\clearpage


\subsection{2015年天津卷（文）}\label{gk:2015年天津卷（文）}


\clearpage


\subsection{2015年浙江卷（理）}\label{gk:2015年浙江卷（理）}


\clearpage


\subsection{2015年浙江卷（文）}\label{gk:2015年浙江卷（文）}


\clearpage

